\documentclass[../../main]{subfiles}

\begin{document}

\section{Lógica Proposicional}
\label{sec:matematica:logica-matematica:logica-proposicional}

\begin{defn}[Conectivo Lógico]
  \label{def:matematica:logica-matematica:logica-proposicional:conectivo-logico}

  Um conectivo lógico é um operador utilizado para construir novas
  proposições a partir de proposições previamente dadas.
\end{defn}

\begin{defn}[Negação]
  \label{def:matematica:logica-matematica:logica-proposicional:negacao}

  Seja (P) uma proposição.

  A negação de (P), denotada por

  [
    \neg P,
  ]

  é verdadeira se, e somente se, (P) é falsa.
\end{defn}

\begin{defn}[Conjunção]
  \label{def:matematica:logica-matematica:logica-proposicional:conjuncao}

  Sejam (P) e (Q) proposições.

  A conjunção

  [
    P \land Q
  ]

  é verdadeira se, e somente se, (P) e (Q) são verdadeiras.
\end{defn}

\begin{defn}[Disjunção]
  \label{def:matematica:logica-matematica:logica-proposicional:disjuncao}

  Sejam (P) e (Q) proposições.

  A disjunção

  [
    P \lor Q
  ]

  é verdadeira se, e somente se, pelo menos uma das proposições é
  verdadeira.
\end{defn}

\begin{defn}[Implicação]
  \label{def:matematica:logica-matematica:logica-proposicional:implicacao}

  Sejam (P) e (Q) proposições.

  A implicação

  [
    P \Rightarrow Q
  ]

  é falsa se, e somente se, (P) é verdadeira e (Q) é falsa.
\end{defn}

\begin{defn}[Equivalência Lógica]
  \label{def:matematica:logica-matematica:logica-proposicional:equivalencia-logica}

  Sejam (P) e (Q) proposições.

  A equivalência lógica

  [
    P \Leftrightarrow Q
  ]

  é verdadeira se, e somente se, (P) e (Q) possuem o mesmo valor de
  verdade.
\end{defn}

\begin{defn}[Tautologia]
  \label{def:matematica:logica-matematica:logica-proposicional:tautologia}

  Uma tautologia é uma proposição composta que é verdadeira para toda
  atribuição possível de valores de verdade.
\end{defn}

\begin{defn}[Contradição]
  \label{def:matematica:logica-matematica:logica-proposicional:contradicao}

  Uma contradição é uma proposição composta que é falsa para toda
  atribuição possível de valores de verdade.
\end{defn}

\begin{defn}[Contingência]
  \label{def:matematica:logica-matematica:logica-proposicional:contingencia}

  Uma contingência é uma proposição composta que não é nem uma tautologia
  nem uma contradição.
\end{defn}

\subsection{Tabelas-Verdade}

\begin{prop}[Tabela-Verdade da Negação]

  [
    \begin{array}{c|c}
      P & \neg P \
      \hline
      V & F \
      F & V
    \end{array}
  ]

\end{prop}

\begin{prop}[Tabela-Verdade da Conjunção]

  [
    \begin{array}{c|c|c}
      P & Q & P \land Q \
      \hline
      V & V & V \
      V & F & F \
      F & V & F \
      F & F & F
    \end{array}
  ]

\end{prop}

\begin{prop}[Tabela-Verdade da Disjunção]

  [
    \begin{array}{c|c|c}
      P & Q & P \lor Q \
      \hline
      V & V & V \
      V & F & V \
      F & V & V \
      F & F & F
    \end{array}
  ]

\end{prop}

\begin{prop}[Tabela-Verdade da Implicação]

  [
    \begin{array}{c|c|c}
      P & Q & P \Rightarrow Q \
      \hline
      V & V & V \
      V & F & F \
      F & V & V \
      F & F & V
    \end{array}
  ]

\end{prop}

\begin{prop}[Tabela-Verdade da Equivalência]

  [
    \begin{array}{c|c|c}
      P & Q & P \Leftrightarrow Q \
      \hline
      V & V & V \
      V & F & F \
      F & V & F \
      F & F & V
    \end{array}
  ]

\end{prop}

\subsection{Leis Fundamentais}

\begin{prop}[Lei da Dupla Negação]
  \label{prop:matematica:logica-matematica:logica-proposicional:dupla-negacao}

  Para toda proposição (P),

  [
    \neg(\neg P)
    \Leftrightarrow
    P.
  ]
\end{prop}

\begin{proof}

  Pela tabela-verdade da negação:

  [
    \begin{array}{c|c|c}
      P & \neg P & \neg(\neg P) \
      \hline
      V & F & V \
      F & V & F
    \end{array}
  ]

  As colunas correspondentes a (P) e (\neg(\neg P)) coincidem.

\end{proof}

\begin{prop}[Primeira Lei de De Morgan]
  \label{prop:matematica:logica-matematica:logica-proposicional:primeira-lei-de-morgan}

  Para quaisquer proposições (P) e (Q),

  [
    \neg(P \land Q)
    \Leftrightarrow
    (\neg P)\lor(\neg Q).
  ]
\end{prop}

\begin{proof}

  Pela tabela-verdade:

  [
    \begin{array}{c|c|c|c}
      P & Q & \neg(P \land Q) &
      (\neg P)\lor(\neg Q) \
      \hline
      V & V & F & F \
      V & F & V & V \
      F & V & V & V \
      F & F & V & V
    \end{array}
  ]

  As duas últimas colunas coincidem.

\end{proof}

\begin{prop}[Segunda Lei de De Morgan]
  \label{prop:matematica:logica-matematica:logica-proposicional:segunda-lei-de-morgan}

  Para quaisquer proposições (P) e (Q),

  [
    \neg(P \lor Q)
    \Leftrightarrow
    (\neg P)\land(\neg Q).
  ]
\end{prop}

\begin{proof}

  Pela tabela-verdade:

  [
    \begin{array}{c|c|c|c}
      P & Q & \neg(P \lor Q) &
      (\neg P)\land(\neg Q) \
      \hline
      V & V & F & F \
      V & F & F & F \
      F & V & F & F \
      F & F & V & V
    \end{array}
  ]

  As duas últimas colunas coincidem.

\end{proof}

\begin{prop}[Comutatividade da Conjunção]

  Para quaisquer proposições (P) e (Q),

  [
    P \land Q
    \Leftrightarrow
    Q \land P.
  ]

\end{prop}

\begin{prop}[Comutatividade da Disjunção]

  Para quaisquer proposições (P) e (Q),

  [
    P \lor Q
    \Leftrightarrow
    Q \lor P.
  ]

\end{prop}

\begin{prop}[Associatividade da Conjunção]

  Para quaisquer proposições (P), (Q) e (R),

  [
    (P \land Q)\land R
    \Leftrightarrow
    P\land(Q\land R).
  ]

\end{prop}

\begin{prop}[Associatividade da Disjunção]

  Para quaisquer proposições (P), (Q) e (R),

  [
    (P \lor Q)\lor R
    \Leftrightarrow
    P\lor(Q\lor R).
  ]

\end{prop}

\begin{prop}[Distributividade da Conjunção sobre a Disjunção]

  Para quaisquer proposições (P), (Q) e (R),

  [
    P\land(Q\lor R)
    \Leftrightarrow
    (P\land Q)\lor(P\land R).
  ]

\end{prop}

\begin{prop}[Distributividade da Disjunção sobre a Conjunção]

  Para quaisquer proposições (P), (Q) e (R),

  [
    P\lor(Q\land R)
    \Leftrightarrow
    (P\lor Q)\land(P\lor R).
  ]

\end{prop}

\end{document}
